\documentclass[aps,prd,twocolumn,superscriptaddress,nofootinbib,10pt]{revtex4-2}

\usepackage{amsmath,amssymb}
\usepackage{graphicx}
\usepackage[colorlinks=true,citecolor=blue,linkcolor=blue,urlcolor=blue]{hyperref}
\usepackage{bm}

\begin{document}

\title{If the lightest neutrino mass is the dark-energy scale:\\
a narrow window for neutrinoless double-beta decay}

\author{Justin Pulford}
\email{pulfordj420@gmail.com}
\affiliation{Unaffiliated}

\date{\today}

\begin{abstract}
The dark-energy density fixes a mass scale $\rho_\Lambda^{1/4} \simeq 2.25\,$meV, of the same
order as the neutrino mass splittings. Numerical proximity is not evidence; we only map
consequences. Taking the sharpest arithmetic reading as a scenario --- $m_1 = 2.25\,$meV with
normal ordering --- and evaluating the standard $m_{\beta\beta}$ map at that single point gives
three contributions $(1.52, 2.67, 1.10)\,$meV, $m_{\beta\beta} \in [0.04, 5.30]\,$meV, and
$\Sigma m_\nu = 61.3\,$meV. The lower edge is an unprotected $1.7\%$ accident: it exists only
because the middle term exceeds the other two by $0.045\,$meV and vanishes entirely for
$m_1 \gtrsim 2.32\,$meV; it has no observational weight. The nearer referee is cosmology: the mass
sum sits just inside current combined limits and will be graded there first. Laboratory
discrimination is secondary. Minimal normal ordering cannot exceed $3.69\,$meV, while this
scenario reaches $5.30\,$meV and occupies the gap between them $32\%$ of the time. One planned
experiment overlaps that gap only at the favourable end of its matrix-element range, and only
$10.8\%$ of phases put $m_{\beta\beta}$ above the corresponding threshold; a detection there would
favour the scenario over minimal ordering, while a detection above $5.30\,$meV excludes it. A null
result constrains nothing, at any sensitivity. The $m_{\beta\beta}$ band itself is textbook; only
the placement at this $m_1$ is at issue.
\end{abstract}

\maketitle

% ---------------------------------------------------------------------------
\section{Introduction}\label{sec:intro}

The energy density of dark energy, expressed as a mass scale, is
%
\begin{equation}
  \rho_\Lambda^{1/4} \simeq 2.25\,\mathrm{meV},
  \label{eq:rho}
\end{equation}
%
which lies within an order of magnitude of $\sqrt{\Delta m^2_{21}} \simeq 8.6\,$meV and two orders
of $\sqrt{\Delta m^2_{31}} \simeq 50\,$meV. That numerical proximity has been remarked on for
decades and has motivated model building --- most directly mass-varying-neutrino scenarios that
tie the dark-energy density to the neutrino sector~\cite{FardonNelsonWeiner2004,PecceiMaVaN2005}.
Proximity of scales is not evidence of a relation. We do not treat it as one, propose no mechanism,
and take no position on those that exist.

What follows is narrower: a single arithmetic scenario and its map. Equation~\eqref{eq:rho} has one
sharp reading,
%
\begin{equation}
  m_1 = \rho_\Lambda^{1/4} = 2.25\,\mathrm{meV}, \qquad \text{normal ordering},
  \label{eq:hyp}
\end{equation}
%
with $m_1$ the lightest neutrino mass. We ask only what Eq.~\eqref{eq:hyp} implies once the
standard oscillation inputs are fixed. No cosmological model enters; nothing computed below
depends on why the equality might hold. Numerical proximity is not evidence; we only map
consequences.

The map from the lightest mass to the effective Majorana mass $m_{\beta\beta}$ is textbook, and the
familiar lobster-shaped allowed region in the $(m_{\rm lightest}, m_{\beta\beta})$ plane appears in
every review of neutrinoless double-beta decay~\cite{DellOroReview2016,AgostiniReview2023}.
Evaluating that map at one particular $m_{\rm lightest}$ is not a new calculation. The only claim is
the placement: just above the hard ceiling of minimal normal ordering, with a mass sum that
cosmology will grade before any laboratory experiment can. That placement, and the (mostly modest)
detection probabilities that follow from it, are what we work out.

Section~\ref{sec:calc} fixes the inputs and computes the window, leading with the mass sum.
Section~\ref{sec:floor} explains why the lower edge is nonzero and why it is an unprotected
accident. Section~\ref{sec:test} sets out what would distinguish the scenario from minimal
ordering and what would refute it. Section~\ref{sec:assume} states the assumptions.

% ---------------------------------------------------------------------------
\section{The effective mass}\label{sec:calc}

For three Majorana neutrinos the amplitude for neutrinoless double-beta decay is controlled by
%
\begin{equation}
  \begin{aligned}
    m_{\beta\beta} &= \Big| \sum_{i=1}^{3} U_{ei}^2\, m_i \Big| \\[2pt]
    &= \Big| c_{12}^2 c_{13}^2 m_1 + s_{12}^2 c_{13}^2 m_2 e^{i\alpha_2}
       + s_{13}^2 m_3 e^{i\alpha_3} \Big|,
  \end{aligned}
  \label{eq:mbb}
\end{equation}
%
with $\alpha_2, \alpha_3$ the Majorana phases, which oscillation experiments do not constrain. We
use $\sin^2\theta_{12} = 0.307$, $\sin^2\theta_{13} = 0.022$,
$\Delta m^2_{21} = 7.42 \times 10^{-5}\,\mathrm{eV}^2$ and
$\Delta m^2_{31} = 2.51 \times 10^{-3}\,\mathrm{eV}^2$, consistent with current global
fits~\cite{NuFIT2020,PDG2024}. With $m_2 = (m_1^2 + \Delta m^2_{21})^{1/2}$ and
$m_3 = (m_1^2 + \Delta m^2_{31})^{1/2}$, Eq.~\eqref{eq:hyp} gives
$m_2 = 8.90\,$meV, $m_3 = 50.15\,$meV, and
%
\begin{equation}
  \Sigma m_\nu = 61.3\,\mathrm{meV}.
  \label{eq:sum}
\end{equation}
%
The three contributions to Eq.~\eqref{eq:mbb} are
%
\begin{equation}
  \big( |U_{e1}^2 m_1|,\ |U_{e2}^2 m_2|,\ |U_{e3}^2 m_3| \big)
  = (1.52,\ 2.67,\ 1.10)\ \mathrm{meV}.
  \label{eq:terms}
\end{equation}
%
Varying the phases over their full range gives
%
\begin{equation}
  \boxed{\ 0.04\,\mathrm{meV} \;\le\; m_{\beta\beta} \;\le\; 5.30\,\mathrm{meV}\ }
  \label{eq:window}
\end{equation}
%
with $\simeq 3.3\,$meV typical. The scenario fixes $m_1$ but says nothing about the Majorana
phases, so position within Eq.~\eqref{eq:window} is unconstrained; that is a genuine limitation and
not a marginal one, as Sec.~\ref{sec:test} makes quantitative.

Cosmology is the nearer referee. The mass sum in Eq.~\eqref{eq:sum} is more immediately testable
than anything in the laboratory programme below. Current bounds from the microwave background
combined with baryon acoustic oscillations have reached $\Sigma m_\nu \lesssim 72\,$meV, with some
combinations pressing lower~\cite{DESI2024}. The value in Eq.~\eqref{eq:sum} sits just inside those
limits; if the bound continues to tighten, cosmology will grade this scenario before any
$0\nu\beta\beta$ experiment can. Laboratory discrimination is secondary and is discussed only after
that fact.

% ---------------------------------------------------------------------------
\section{Why the lower edge is nonzero, and why it is fragile}\label{sec:floor}

The lower edge of Eq.~\eqref{eq:window} is not a robust floor; it is an unprotected numerical
accident, and that status should be stated before any arithmetic. Complete cancellation in
Eq.~\eqref{eq:mbb} requires the three complex terms to close a triangle, which is possible only if
no one of them exceeds the sum of the other two. From Eq.~\eqref{eq:terms},
%
\begin{equation}
  2.67 \;>\; 1.52 + 1.10 = 2.62 \ \mathrm{meV},
  \label{eq:triangle}
\end{equation}
%
by $0.045\,$meV. The middle term wins, the triangle cannot close, and the lower edge follows as the
residue $2.67 - 2.62$.

That margin is only $0.045\,$meV on terms of order $2.6\,$meV --- a $1.7\%$ effect. Nothing
protects it: it is not the consequence of a symmetry, and it would not survive a modest shift in
$\theta_{12}$, in $\Delta m^2_{21}$, or in $m_1$ itself. At $m_1 \gtrsim 2.32\,$meV the inequality
reverses and the floor vanishes entirely. A $3\%$ change in the lightest mass, or a comparably small
shift in the oscillation inputs, erases the nonzero lower edge. Calling it a floor overstates what
it is: an unprotected numerical accident at this particular $m_1$.

The same near-degeneracy makes $m_{\beta\beta}$ a steep function of $m_1$ near this point. A quantity
sitting close to a cancellation responds sharply to its inputs, so the band edge here is unusually
sensitive to the inputs that produce it.

No observational consequence rides on the lower edge. At $0.04\,$meV it is two orders of magnitude
below any experimental reach under discussion. Every statement in Sec.~\ref{sec:test} concerns the
upper edge, which the near-cancellation barely moves.

% ---------------------------------------------------------------------------
\section{What would test it}\label{sec:test}

\subsection{The interval that separates scenarios}

Laboratory tests are secondary to the cosmological mass sum of Sec.~\ref{sec:calc}. Conditional on
that, the comparison that matters in $0\nu\beta\beta$ is not against zero but against the minimal
alternative. Setting $m_1 = 0$ with the same ordering and parameters gives
%
\begin{equation}
  m_{\beta\beta}^{\rm min} \in [1.48,\ 3.69]\ \mathrm{meV},
  \label{eq:minimal}
\end{equation}
%
so minimal normal ordering has a hard ceiling at $3.69\,$meV that no choice of phases can exceed.
The scenario of Eq.~\eqref{eq:hyp} reaches $5.30\,$meV. The interval
%
\begin{equation}
  3.69\,\mathrm{meV} \;<\; m_{\beta\beta} \;<\; 5.30\,\mathrm{meV}
  \label{eq:band}
\end{equation}
%
is therefore inaccessible to minimal ordering and occupied by this scenario. Integrating
Eq.~\eqref{eq:mbb} over flat Majorana phases, the scenario lands inside Eq.~\eqref{eq:band}
$31.8\%$ of the time.

\subsection{Where the experiments fall}

Projected ten-year reaches for the ton-scale programme, quoted as ranges because the nuclear matrix
element is the dominant uncertainty, are approximately $4.7$--$20.3\,$meV for
nEXO~\cite{nEXO2021}, $9$--$21\,$meV for LEGEND-1000~\cite{LEGEND1000}, and $12$--$34\,$meV for
CUPID~\cite{CUPID2019}.

Against the ceiling of Eq.~\eqref{eq:window}, LEGEND-1000 and CUPID lie entirely above the
scenario and cannot reach it for any matrix element. Only nEXO overlaps, and only at the favourable
end of its matrix-element range, $4.7$--$5.3\,$meV. Over flat phases the scenario exceeds
$4.7\,$meV only $10.8\%$ of the time --- roughly one in nine. Even if Eq.~\eqref{eq:hyp} is
correct, baseline nEXO is unlikely to see a signal: low probability is the first statement about
this overlap, not a footnote to it.

Discrimination is secondary to that probability. The entire $4.7$--$5.3\,$meV overlap does lie
inside Eq.~\eqref{eq:band}, which minimal ordering cannot reach, so a detection there would favour
the scenario over minimal normal ordering. It would not be uniquely attributable to
Eq.~\eqref{eq:hyp}: nuclear-matrix-element uncertainty, phase measure, and the secondary status of
the laboratory test relative to cosmology all limit how much weight such a signal could carry.

\subsection{Why a sensitivity upgrade does not help as much as expected}

A factor-of-four improvement in half-life sensitivity, of the kind a barium-tagging upgrade is
projected to deliver, corresponds to a factor $\sqrt{4} = 2$ in $m_{\beta\beta}$, taking the reach
to $\simeq 2.35\,$meV. The detection probability rises from $10.8\%$ to $69.1\%$.

This buys likelihood, not discrimination. Minimal normal ordering also exceeds $2.35\,$meV, and
does so $63.7\%$ of the time. A tagged detection near that threshold would be consistent with both
scenarios and would separate almost nothing. The unimproved machine is the more discriminating one
in $0\nu\beta\beta$, precisely because its reach sits in the interval where the two scenarios do not
overlap --- still subject to the low $10.8\%$ probability above, and still secondary to the
cosmological mass sum.

\subsection{Falsification}

A measured $m_{\beta\beta} > 5.30\,$meV refutes Eq.~\eqref{eq:hyp} outright, since no choice of
phases can produce it. Evidence for inverted ordering does the same: at $m_3 = 2.25\,$meV inverted
ordering gives $m_{\beta\beta} \in [18.6, 49.1]\,$meV, which does not overlap
Eq.~\eqref{eq:window} at any phase. Evidence that neutrinos are Dirac particles removes the
observable entirely. Cosmological exclusion of $\Sigma m_\nu = 61.3\,$meV would do the same, and
is the nearer path: the laboratory ceiling is a secondary falsifier.

A null result in $0\nu\beta\beta$ refutes nothing. Because the phases are free, $m_{\beta\beta}$ can
be arbitrarily small, and no non-observation at any laboratory sensitivity constrains
Eq.~\eqref{eq:hyp}. That asymmetry is built into the scenario.

For orientation, the strongest current limit comes from $^{136}$Xe: KamLAND-Zen reports
$T_{1/2} > 4.3 \times 10^{26}\,$yr at $90\%$ confidence over a $2097\,$kg\,yr exposure,
corresponding to $m_{\beta\beta} < 28$--$122\,$meV across the matrix-element
range~\cite{KamLANDZen2023}. That is a factor of five or more above the ceiling of
Eq.~\eqref{eq:window}, so no existing laboratory measurement bears on this scenario in either
direction.

% ---------------------------------------------------------------------------
\section{What is assumed}\label{sec:assume}

\emph{Majorana neutrinos.} Eq.~\eqref{eq:mbb} presumes lepton-number violation. If neutrinos are
Dirac there is no $m_{\beta\beta}$ and the question does not arise.

\emph{Normal ordering.} Assumed throughout, and stated in Eq.~\eqref{eq:hyp} rather than derived.
Global fits currently prefer it at modest significance; the preference is not settled.

\emph{Three light species.} A light sterile state would add a term to Eq.~\eqref{eq:mbb} and widen
the window.

\emph{Standard mechanism.} We assume light-neutrino exchange dominates the amplitude. If other
lepton-number-violating physics contributes, $m_{\beta\beta}$ as extracted from a half-life is no
longer the quantity computed here.

\emph{The numerical input.} We use the observed $\rho_\Lambda^{1/4} = 2.25\,$meV. Nothing depends
on this at the quoted precision: shifting it by $0.5\%$ moves the lower edge of
Eq.~\eqref{eq:window} from $0.045$ to $0.039\,$meV and leaves the upper edge at $5.30\,$meV, so
every conclusion in Sec.~\ref{sec:test} is unchanged.

\emph{The scenario itself.} Eq.~\eqref{eq:hyp} is an arithmetic scenario, not a result. Equating
two scales because they are numerically close is not evidence; we do not argue otherwise. We write
it down because it is sharp, costs nothing to state, and happens to be nearly testable --- first by
cosmology through $\Sigma m_\nu$, secondarily in the laboratory. Numerical proximity is not
evidence; we only map consequences.

% ---------------------------------------------------------------------------
\section{Conclusion}

Taking $m_1 = 2.25\,$meV with normal ordering as a sharp arithmetic scenario --- not as evidence
drawn from the coincidence of scales --- gives $m_{\beta\beta} \in [0.04, 5.30]\,$meV and
$\Sigma m_\nu = 61.3\,$meV. The lower edge is an unprotected $1.7\%$ accident ($0.045\,$meV margin;
vanishes for $m_1 \gtrsim 2.32\,$meV) and carries no observational weight. The upper edge does the
laboratory work; the mass sum is the nearer number.

Cosmology is the first referee. A sum of $61.3\,$meV sits just inside current combined limits, and
the bound is tightening. Laboratory discrimination is secondary. Conditional on that, the useful
$0\nu\beta\beta$ comparison is against minimal normal ordering, whose hard ceiling is $3.69\,$meV.
The gap between the two ceilings is closed to minimal ordering and occupied by this scenario about
a third of the time. One planned experiment overlaps that gap only at favourable matrix elements,
and only $10.8\%$ of phases put a signal above the corresponding threshold: low probability first,
modest discrimination second. A detection there would favour the scenario over minimal ordering
without uniquely establishing it; a detection above $5.30\,$meV would refute the scenario; a null
laboratory result would say nothing at all.

\bibliographystyle{apsrev4-2}
\bibliography{refs}

\end{document}
